% ===== §8 The appropriate poem: the non-possessive interpretive cycle that generates real value =====
\section{The Appropriate Poem: The Non-Possessive Interpretive Cycle That Generates Real Value}
\label{sec:core}

\noindent
Up to here the greater part of the paper has been spent diagnosing the vicious---failure, extraction, the stolen phase, the outside reduced to fuel. This was a necessary detour, but the time has come to put the question affirmatively: since poeticity does not guarantee the good, \emb{what is the appropriate, the just, the self-perpetuating poeticity}?

This section is the keystone of the paper. Its particularity is this: the answer is not derived along a single line, but is the meeting, at one and the same point, of \emph{two independent inquiries, each walked to its end}. One asks ``what is the appropriate poeticity,'' setting out from form and ethics; the other asks ``how is value generated in an intimate relation,'' setting out from a theory of value. They converge on a single thesis. And this convergence itself---two different frameworks illuminating one another, climbing in a spiral, carrying home a remainder neither foresaw alone---is a live occurrence of the very good circle the paper asserts. The form of this section, once again, enacts its content.

\subsection{The first approach: the appropriate poem is the non-possessive poem}
\label{sec:non-possess}

We cannot define ``the good poeticity'' as ``poeticity $+$ the good'' and declare ourselves finished---that is tautology. The genuinely interesting question is: does the vicious poetic cycle, \emph{within} its poetic structure, leave a recognizable formal scar? That is---does ``stealing the phase by extraction'' leave a mark upon the manner in which the poem operates, such that we might smell the vampire from \emph{within} poeticity, rather than by appending a criterion of justice from without?

I claim: it does. For ``to reduce some other to pure fuel'' is a \emph{closing} operation upon interpretation, and it necessarily leaves three scars upon the interpretive structure of the poem. And the obverse of these three scars sketches, exactly, the contour of the ``appropriate poem.'' More remarkable still---they fall in three distinct frameworks, mutually irreducible, and yet converge upon a single word.

\subsubsection{The hermeneutic dimension: can it be freely betrayed}

Return to the poem's root act: the poem's generativity unfolds in the interpreter, who participates in person in the generation of meaning. The true poem installs a living grammar the interpreter may \emph{freely derive from}---the sentences you generate with it are \emph{yours}, bearing your path, your situation, your symmetry-breaking this once. The poem does not prescribe where you must arrive.

The poeticity of the demagogue---fascist aesthetics, cultic speech---installs, by contrast, a grammar with \emb{but one legitimate derivation}. On the surface its $\lambda>1$ (it ceaselessly generates new content, new passion), yet all derivations converge, covertly, upon one predetermined endpoint: the leader, the enemy, salvation. What it demands is not witness but \emph{submission}. The interpreter believes himself to be generating freely, while in fact executing a conclusion long preinstalled.

\begin{claim}[The first scar: possessing the reading]
This is the disguised version of settling failure: it does not publicly declare the loop closed (that would expose it), but lets you believe the loop open while covertly steering every path to one closing point. Hence the first criterion of the appropriate poem is hermeneutic: \emph{its grammar permits genuinely branching derivation, handing to the interpreter the generativity together with the freedom to go elsewhere, even to oppose it---it can be freely betrayed.}\footnote{This is the hermeneutic insight that understanding is never the passive reception of a fixed meaning but a productive event in which interpreter and text co-determine sense \citep{gadamer2004}; the genuine work outruns any single reading, a \emph{surplus of meaning} that no interpretation exhausts \citep{ricoeur1976}.} A true poem permits, even invites, a reading that betrays it; a vicious poetic grammar permits no betrayal---to betray is to apostatize, to become the enemy. Whether it can be freely opposed is the watershed between the poem and the demagogue.
\end{claim}

\subsubsection{The psychoanalytic dimension: encircling the void, or the idol}

The second scar continues from the bearer fixed in the preceding section---the sublimation of jouissance. The structure of sublimation is this \citep{lacan1992}: desire \emph{encircles} an essentially empty, impossible real (\emph{das Ding}) and does not fill it; and just because it does not fill it, desire can go on generating, enduring, sublimating. \emb{The true poem keeps the void as void}---it points toward the unsayable without pretending to say it. This is the source of its inexhaustibility: that central emptiness guarantees that no single interpretation can fill it, settle it.

The vicious poem does the contrary: it \emb{fills the void of \emph{das Ding} with a positive object}---the leader, the race, purity, the enemy. This lends it a vast and immediate energy (a diffuse anxiety quieted by one concrete answer), but at a fatal cost: once the void is filled, that filler becomes an idol to be defended and fed without end, and \emph{every other who threatens the filler is reduced to fuel to be cleared away}.

\begin{claim}[The second scar: possessing the void]
This is the psychoanalytic mechanism of the ``stolen phase'': the vicious poem fills with a sacrifice the void that ought to stay open, and confirms the filler over and over by clearing away others. \emph{The one who encircles an idol needs an enemy} (the idol is forever threatened); \emph{the one who encircles a void needs no sacrifice} (the void cannot be threatened). Hence the second criterion of the appropriate poem is psychoanalytic: it encircles a void kept open rather than filling it, and so its sublimation is \emph{objectless} (encircling the void, self-sufficient, extracting no outside). This is the precise form, in the psychoanalytic framework, of the difference between ``phase internal'' and ``phase stolen.''
\end{claim}

\subsubsection{The political-economic dimension: phase refluxing, or siphoned}

The third scar is political-economic. The true poem is a \emb{decentred} generativity---each interpreter generates, along his own path, a sublimation that is his own, and the phase \emph{stays in the hands of the one who generates it}. The poem does not possess its readings, nor do readers contend over a scarce ``correct reading.'' It is positive-sum: your reading does not diminish mine but may illuminate it.

The vicious poem is a \emb{siphoning} generativity---the whole of the meaning, passion, and loyalty the interpreters generate is drawn back to a centre (the leader, the organization, the brand). The follower witnesses in person, sublimates in person, but the \emph{proceeds of the sublimation are alienated}---they flow to the centre and do not remain in the follower's own hands. The follower is the producer of the phase and yet the one dispossessed of it.

\begin{claim}[The third scar: possessing the proceeds of sublimation]
This is the \emph{alienated} poem: the interpreter contributes the whole labour of witness, yet does not own the fruit of his witnessing. The cult, the fascism, the fanatic brand---each has a centre that draws all in; the true poem has none---an old poem does not take your reading into the author's keeping. Hence the third criterion of the appropriate poem is political-economic: it is decentred, the phase refluxing to every witness, with no centre siphoning the whole of the sublimation. This is the form, in the political-economic framework, of ``no one reduced to pure fuel''---for in the siphoning structure the follower is himself the consumed fuel, though he take himself for a beneficiary.
\end{claim}

\subsubsection{The four frameworks converge on ``non-possession''}

Joining the three scars, the contour of the appropriate poeticity grows clear, and it falls beautifully across three mutually irreducible frameworks---which is no coincidence but a further endogenous corroboration of the paper's polyphonic structure. The three scars are, in truth, one and the same thing showing itself in three frameworks: the vicious poem \emb{possesses} the interpreter's path (permits no betrayal), \emb{possesses} the void (fills it with an idol), \emb{possesses} the proceeds of sublimation (siphons them to a centre).

And \emph{possession} is just another name for \emph{extraction}. Therefore:

\begin{claim}[The appropriate poem is the non-possessive poem]
A poeticity is appropriate (good, sustainable, just) if and only if it \emph{does not possess}: hermeneutically, its living grammar permits genuinely branching derivation and can be freely betrayed; psychoanalytically, it encircles an open void rather than filling it, and so needs no sacrifice; politically-economically, it is decentred, the phase refluxing to every witness rather than siphoned to a centre. These three ``non-possessions'' are three mutually irreducible faces of a single disposition.
\end{claim}

These three non-possessions receive, in the fourth framework---the Daoist---their unifying name. ``To generate and foster them; to generate without possessing, to act without presuming, to foster without ruling: this is called the dark virtue'' (\zh{生而不有，为而不恃，长而不宰，是谓玄德}). To generate without possessing what one generates, to act without leaning on what one does, to foster without lording over---this is, almost word for word, the definition of the appropriate poeticity. The vicious poem is, conversely, the poeticity of ``possessing, presuming, ruling'': it possesses the reading, presumes upon the passion, rules over the sublimation.

So the four frameworks all converge, on this question, upon a single point---``non-possession''---and yet each speaks in its own tongue, none subsumed by another: hermeneutics says ``permits betrayal,'' psychoanalysis says ``encircles the void and not the idol,'' political economy says ``the phase refluxes and is not siphoned,'' Daoism says ``generates without possessing.'' They are four mutually irreducible faces of one ``non-possession.'' \emb{And this, in itself, is a live demonstration of a good poetic cycle}: four voices encircle one void (``the appropriate poem,'' a centre fillable by no single framework), each sublimating, the phase refluxing to every framework, none reduced to a footnote of another.

\subsection{The second approach: how value is generated in an intimate relation}
\label{sec:value}

Now enter by a wholly different road. Throughout we have spoken of ``phase,'' ``sublimation,'' ``holonomy,'' yet have all along treated \emph{value} as an unanalyzed primitive---we said ``the phase internal or stolen,'' ``the sublimation refluxing to every witness,'' but the sublimation of \emph{what}? the reflux of \emph{what}? Without opening this black box of ``value,'' every criterion above hangs in the air. This subsection opens it.\footnote{The scope must first be declared: the task of this subsection is single and precise---to identify, for the ``phase / value'' left hanging above, a stratified bearer, so that the criteria acquire a concrete referent. It does \emph{not} undertake to build here a complete theory of value; that is the work of a separate study (which the series intends to take up in a sequel). Here we use it and pass on, lest the theory of value usurp the stage and drown the true theme of sustainability.}

I claim that value in an intimate relation falls into three strata, answering exactly to the three Lacanian registers, and answering, too, exactly to the structure of good and vicious poeticity we have already drawn out.

\emb{Imaginary value} comes from the imaginary: idealizing the other into a complete image, imagining the relation as two made one, each filling the other's lack. It is the fuel of passionate love, the thing the rose is at its most resplendent. But it is, structurally, that operation of ``filling the void of \emph{das Ding} with a positive object''---it takes the other as the object that could fill my lack (mistaking the \emph{objet petit a} for a possessable positive thing) \citep{lacan2006}. Thus imaginary value tends, by nature, toward the possessive, idol-encircling poeticity: the phase it manufactures is intense yet stolen---idealization is necessarily attended by the erasure of the real other (you love the other of your imagining, while the real other is reduced to fuel for that image). \emb{The rose blooms but a week precisely because imaginary value is exhaustive}: the intensity of idealization cannot be sustained, and the phantasm is, in the end, pierced by the real other. Imaginary value is the psychoanalytic name of ``resplendence born of unsustainability.''

\emb{Symbolic value} comes from the symbolic: status, money, the recognized identity, the displayable sign of relation. It is countable, commensurable, accumulable---and therefore \emph{stealable, siphonable, inflatable, cheaply machine-generable}. It exists, undeniably (to deny the office of the sign is naive); but its danger is just the danger of the sign. When a relation's value-cycle runs in the symbolic stratum alone, it falls into settling or inflationary failure. Symbolic value is the bearer in which the phase is \emph{most easily stolen}, for it is commensurable and extractable---it is the stratum on which the vicious cycle most readily lodges.

\emb{Real value} comes from the real: the stratum bound up with trauma, with jouissance, with the unsymbolizable. It cannot be settled (it commensurates into no sign), cannot be idealized (it is gathered into no complete image of the imaginary). It is precisely that void of \emph{das Ding} itself---the abyss between two persons that can never be wholly mutually understood, never made one, that keeps an irreducible alterity.

\begin{claim}[The three strata of value and the good/vicious poem]
\emph{Imaginary value} $=$ the idol-encircling poem $=$ the possessive sublimation $=$ the stolen phase (idealization erasing the real other) $=$ the resplendence and the swift death of the rose. \emph{Symbolic value} $=$ the settleable, siphonable poem $=$ the most easily extracted bearer $=$ the hotbed of inflationary and settling failure. \emph{Real value} $=$ the void-encircling poem $=$ the non-possessive sublimation $=$ the internal phase $=$ the true bearer of the good's sustainable circle.
\end{claim}

The bond of real value to trauma must be drawn with care, lest one fall into a ``cult of trauma''---as though the more traumatic, the more unsayable, the more valuable. This is false. In a real relation, real value includes also the calm, non-traumatic unsayable: the shared silence, the bodily presence, the daily mutual witness. So it should be said: \emb{real value has trauma as its sharpest, least evadable manifestation, but is not confined to trauma}---it is every relational kernel unsettleable by sign and unpossessable by the imaginary, of which trauma is merely the most compelling form. In an intimate relation the truly deep connection is built, more often than not, upon \emph{witnessing the other's irreducible void without trying to fill it / repair it / possess it}---you accompany the other's unfillable void (trauma, finitude, the fact of eventual parting) rather than cover it with the idealization of the imaginary or settle it with a solution of the symbolic. This, precisely, is the relational form of ``encircling the void, not the idol.''

\subsubsection{Why value must be generated by hermeneutic engagement}

From this a key question can be answered: why \emph{must} real value be generated by hermeneutic engagement? Because real value cannot be given directly (it is unsymbolizable); it can only be approached \emph{encirclingly}, along the path of interpretation and witness. You cannot ``possess'' the other's real kernel; you can only interpret it, witness it, accompany it, time and again---and this just \emph{is} the poetic interpretive cycle.

\begin{claim}[The generative mechanism of real value is the poetic interpretive mechanism]
Real value is not a conserved quantity transmitted; it is a phase \emph{generated} in the process of two persons interpreting and witnessing, again and again, that unfillable void. It does not pre-exist, awaiting discovery; it is generated in the process of being encircled, of being witnessed. Hermeneutic engagement is necessary precisely because the good's value (real value) is, ontologically, interpretively generated---without the path of interpretation there is no such value. Thus \emb{the generative mechanism of real value just is the poetic interpretive mechanism.}
\end{claim}

\subsection{The convergence of the two approaches}
\label{sec:convergence}

Now the two roads meet. The first asked ``what is the appropriate poeticity,'' and answered: the interpretive cycle unfolded \emph{non-possessively} (encircling the void, permitting betrayal, the phase refluxing). The second asked ``how is the good value of an intimate relation generated,'' and answered: real value---that value alone unpossessable, unsiphonable, approachable only by encircling witness---is generated through the poetic interpretive cycle. The two answers are one:

\begin{claim}[The keystone thesis]
\emph{The appropriate poem $=$ the non-possessive interpretive cycle that generates real value.} It is, in form, poetic (a living grammar, enduring, resistant to settling); in justice, good (non-possessive $=$ non-extractive $=$ phase internal $=$ no sacrifice $=$ decentred and refluxing). It neither settles (unlike the contract), nor siphons (unlike the demagogue), nor idles (unlike the false poem). It generates without possessing what it generates---and therefore it can perpetuate itself, precisely \emph{because} it does not try to possess its own perpetuation. This is the final form of ``let the circulation of value become a good, natural circle, and it will perpetuate itself of itself'': \emb{the non-possessive generativity perpetuates itself.}
\end{claim}

Note how this thesis returns upon the paper's point of departure. The metaphysical ground laid in \xref{sec:thesis} was ``generativity takes its own continuation as its end''; now the terminus (non-possession) returns upon the origin (self-continuation)---\emph{it is precisely because it does not possess that self-continuation becomes possible}. To possess is to want to clutch the cycle's product, to clutch the balance of the sublimation, to clutch the continuation itself---and the clutching is precisely what severs the cycle (this is another saying of \xref{sec:praxis}'s ``the one who acts upon it ruins it, the one who grasps it loses it''). Non-possession is not the price of perpetuation but its condition. The distance between terminus and origin is itself a spiral return---a return to ``self-continuation,'' carrying home the remainder ``non-possession'' that was not there at the outset. \emb{The paper has, in its form, completed the good circle it asserts.}

Here, at last, the bearer of the phase can be finally fixed: it is \emb{the sublimation of real value}---the jouissance generated, encirclingly, in each witness through the poetic interpretive cycle, and refluxing to each witness. This at once discharges the two debts left by \xref{sec:synthesis} and \xref{sec:semiotics}: the phase has a real bearer (no longer hanging in the air), and ``sublimation is positive phase'' has a content independent of definition (it is the remainder internal to the non-possessive cycle, as against the counterfeit phase stolen by the possessive cycle).

\subsection{The coupled system as touchstone: returning value to every co-creator}
\label{sec:coupled}

Finally, all of this must be pushed from the purely dyadic relation into its real social embedding. An intimate relation is never an isolated two-person system; it is \emph{coupled} into a larger system---family, kin, society. And the criterion of the good circle is tested most severely precisely within this coupled system.\footnote{This subsection unfolds throughout in generalized language (``coupled system,'' ``participant,'' ``co-creator''); the criteria it states apply to any intimate system embedded in a larger relational network, and depend on no particular cultural form or case.}

Here the hardest principle of generative justice steps forward: \emb{value should return to the creator of the value himself.} Carried into the intimate coupled system, it at once illuminates a fact long obscured by romantic talk---that an intimate relation is a system of the \emph{co-production} of value, and that this system is exceedingly prone to alienation. Who creates the value? Not the loving two alone, but the family members drawn in: they contribute care, labour, affection, symbolic recognition, even material support. All these are the creation of value. Hence the criterion of the good demands:

\begin{claim}[The good circle of the coupled system]
An intimate coupled system has a good value-cycle if and only if: (1) value \emph{can circulate}---guaranteed by the poetic language: value is not settled, not siphoned, but continually reproduced in the loop of interpretation and witness; and (2) value \emph{returns to every one who creates it}---this is the criterion of generative justice: whoever takes part in the creation of value ought to benefit from the value the relation generates, and in the system there exists \emph{no ``pure fuel,'' no one who only creates and never benefits}. The first condition is dynamical (does the cycle live, is it a spiral); the second is just (on whose blood does the cycle run, to whom does it return).
\end{claim}

This is the confluence, here, of the orthogonal structure ``poeticity $\times$ justice'' laid down in \xref{sec:turn}: the good circle must satisfy \emph{both} at once. And the two are not two parallel checklists---they need one another. This is the original kernel of this subsection:

\begin{claim}[Poeticity is the mechanism that makes return possible; return is the content of poeticity's goodness]
The poetic language is the guarantor of the good circle precisely because it is the only form of cycle that lets value ``return to the creator himself'' without alienation: in the poetic cycle, value (the sublimation of real value) is generated along each witness's own path and stays in his own hands. By contrast the purely symbolic value-cycle tends, by nature, to siphon---someone is always serving as fuel for another's symbolic capital. Therefore only when the value-cycle is led by \emph{real value} and runs in the poetic manner does ``the creator is the beneficiary'' become structurally possible; for real value cannot be possessed, cannot be siphoned into another's name. Poeticity and return are two faces of one thing: poeticity is the mechanism that makes return possible, return is the content of poeticity's being good.
\end{claim}

From which it appears that this section makes a double advance upon generative justice. The first is the \emb{temporal dimension} (the axiom of sustainability of \xref{sec:synthesis}): the return must not only occur but suffice to reproduce the conditions of generation themselves. The second is the \emb{problem of the bearer of return}: heretofore generative justice has mostly spoken of return in material, labour, ecological value, but in the intimate domain the least alienable value is real value, and it \emph{can} return only through the poetic interpretive cycle. Generative justice in the intimate domain is therefore necessarily \emph{poetic}---in a structural and not a metaphorical sense: only the non-possessive, non-settling, phase-refluxing cycle can return that unsymbolizable value to every one who creates it. \emb{The intimate relation is thus the limit case that tests whether generative justice can handle ``the incommensurable, unpossessable value,''} and the poetic cycle is its realization in that limit case.

At the scale of the family, alienation has a typical form---the alienation of social reproduction. The labour of care, of affect, of the household creates the value that sustains the relation, yet because it is invisible, because it is idealized as ``the natural outflowing of love,'' it is not returned: one party goes on creating value yet does not benefit from the value she creates, reduced to pure fuel within this system of tenderness.\footnote{This is the structure disclosed by social-reproduction theory \citep{federici2004,bhattacharya2017}: the reproductive labour that sustains life and labour-power creates value, yet is placed, by the twin narratives of capital and patriarchy, in the position of the invisible and the uncounted. Fraser names the resulting instability a \emph{crisis of care}---capital depends on reproductive labour it refuses to value, and so erodes the very condition of its own continuation \citep{fraser2017}, which is, in the vocabulary of \xref{sec:antithesis}, the counterfeit closure seen at the scale of the household.} This is just what ``the idol-encircling poem'' looks like within the family: with idealized signs---``virtue,'' ``sacrifice,'' ``devotion''---it whitewashes one who is extracted from into one who gives of her own accord; \emph{imaginary value here performs the cover for extraction}.

The coupled system, by reason of its internal power differentials, its external pressures, its conflict of multiple value-systems (the family is often the executor of symbolic value), makes ``no one is pure fuel'' pass from an abstract proposition into a real demand to be vigilantly, ceaselessly diagnosed. So the practice of the good takes, within the coupled system, a concrete form: to let the rite \emph{anchor without settling} real value (the rite as an invitation to the path, not the discharge of a debt), and to interrogate, without end, the question of the pure-fuel criterion (\xref{sec:fuel})---\emb{is there, in this enlarged cycle, anyone who creates value yet does not benefit}?

And this returns us to a beautiful gathering. The essence of the good circle is \emb{co-authorship}---which is precisely the unfolding, in the multi-agent system, of the hermeneutic dimension of \xref{sec:non-possess}: ``permitting free betrayal, handing the generativity truly to the interpreter.'' In the good coupled system every participant is a \emph{co-author} of the relation's meaning, not the executor of a written role: none is assigned a script of ``virtue'' to perform, none is the passive object receiving a settled arrangement; each can possess, in the generation of the relation's meaning, his own branching derivation, his own path, his own phase. \emb{To co-create value is to co-author a poem possessed by no one.}
